% -----------------------------------------------------------------------------
% 4. Results — un résultat DÉMONTRÉ (coût du format vs effet de granularité),
% deux résultats RAPPORTÉS (plafond humain / juges ; modèle fine-tuné).
% Budget : ~450 mots + Tables 1 et 2.
% -----------------------------------------------------------------------------
% Table 1 — le résultat démontré. Légende AU-DESSUS (règle IOS Press).
\begin{table}[t]
\caption{Reliability and downstream signal by taxonomy, on all 50 contracts
(9,414 sentences, three annotators). The gap between $\alpha$-MASI and nominal
$\alpha$ is the cost of the multi-label format. Differences against \Tn{20} are
paired by document; 95\% confidence intervals from 1{,}000 document-level
bootstrap resamples. Signal loss is the relative drop in leave-one-document-out
average precision of $P(\mathrm{unfair}\mid\mathrm{theme\ combination})$, scored
against the \corpus{} unfairness labels, which are never projected.
\Tn{10} merges across unfairness strata; \Tn{14} and \Tn{11} do not.}
\label{tab:reliability}
\centering
\begin{tabular}{lccccc}
\toprule
Taxonomy & Classes & $\alpha$-MASI & Nominal $\alpha$ & $\Delta$ vs.\ \Tn{20} [95\% CI] & Signal loss \\
\midrule
\Tn{20} & 20 & 0.658 & 0.732 & --- & --- \\
\Tn{14} & 14 & 0.693 & 0.768 & $+0.035$ [0.030, 0.042] & 6.4\% \\
\textbf{\Tn{11}} & 11 & \textbf{0.725} & 0.794 & $\mathbf{+0.067}$ \textbf{[0.058, 0.077]} & \textbf{5.9\%} \\
\Tn{10} & 10 & 0.737 & 0.800 & $+0.078$ [0.068, 0.090] & 13.2\% \\
\bottomrule
\end{tabular}
\end{table}


\paragraph{The cost of the multi-label format.}
Scoring secondary themes lowers agreement from a nominal $\alpha$ of 0.732 to an
$\alpha$-MASI of 0.658: the multi-label format costs $0.073$ [0.061, 0.087] of
chance-corrected agreement. Across the four granularities it stays between 0.063 and
0.075, without decreasing monotonically as categories are merged
(Table~\ref{tab:reliability}); its whole range, 0.012, is five times smaller than
the agreement gain consolidation yields below.

\paragraph{The effect of granularity.}
Consolidating the taxonomy recovers agreement of the same order as the format
costs. Merging \Tn{20} into \Tn{11} raises $\alpha$-MASI by $0.067$ [0.058, 0.077],
carrying the layer across the 0.667 threshold conventionally used for tentative
conclusions \cite{passonneau2006masi} --- a threshold the annotation taxonomy
itself does not reach. Two properties make this more than a mechanical artefact
of coarser categories. First, the gain \emph{replicates out of sample}: the
17 contracts held out from the design of the merges yield $+0.069$, against
$+0.066$ on the 33 contracts used to design them. Second, \Tn{11} \emph{dominates}
the intermediate \Tn{14}: it gains twice as much agreement ($+0.067$ against
$+0.035$) while losing \emph{less} downstream signal (5.9\% against 6.4\%).
Granularity is therefore a genuine lever on reliability, and not merely a
redistribution of the same disagreements over fewer categories.

\paragraph{The legal constraint on consolidation.}
Consolidation cannot be driven by agreement alone. \Tn{10}, obtained by merging
along measured confusion clusters without regard to the legal status of the
clauses, buys a further $0.012$ of $\alpha$-MASI over \Tn{11} --- and loses
13.2\% of the unfairness signal instead of 5.9\%, a yield of 0.59 points of
agreement per unit of signal against 1.13 for \Tn{11}. The merge responsible joins two
themes whose observed unfairness rates fall in opposite strata, 0.385 against
0.039 --- a distinction consumer protection law draws
\cite{eu1993directive,micklitz2017empire} and that the agreement signal does not
express. A purely statistical consolidation selects exactly this taxonomy;
the legal constraint rules it out: reliability gains and downstream utility are
not interchangeable, and the boundary between them is drawn by law, not by data.

% Table 2 — les résultats rapportés. Une seule référence de comparaison ; les
% quatre juges sont nommés et classés (E3.3 : classement identique en T20 et T11).
\begin{table}[t]
\caption{Agreement with the frozen gold standard under the annotated (\Tn{20}) and
the consolidated (\Tn{11}) taxonomy. Human agreement is leave-one-annotator-out, each
annotator being scored against a gold to which they did not contribute; the encoder is
evaluated 5-fold by document (95\% intervals from 1{,}000 document-level bootstrap
resamples of its out-of-fold predictions); judge labels were produced once and stored.
Judges are ordered by agreement, identically under both taxonomies. Reported, not
demonstrated: see Section~\ref{sec:results}.}
\label{tab:systems}
\centering
\begin{tabular}{lcccc}
\toprule
 & \multicolumn{2}{c}{$\kappa$} & \multicolumn{2}{c}{Accuracy} \\
System & \Tn{20} & \Tn{11} & \Tn{20} & \Tn{11} \\
\midrule
Human annotators & \textbf{0.859} & --- & 0.871 & --- \\
Legal-BERT, fine-tuned & 0.695 [0.671, 0.720] & 0.720 [0.692, 0.748] & 0.719 & 0.754 \\
Judge \textsc{Fable} (Claude Fable 5) & 0.581 & 0.591 & 0.601 & 0.629 \\
Judge \textsc{Claude} (Claude Opus 4.7) & 0.515 & 0.531 & 0.539 & 0.576 \\
Judge \textsc{Mistral} (Mistral Medium 3.5) & 0.327 & 0.353 & 0.358 & 0.413 \\
Judge \textsc{Codex} (GPT-5.5) & 0.266 & 0.302 & 0.305 & 0.375 \\
\bottomrule
\end{tabular}
\end{table}


\paragraph{Reported: the human ceiling, the judges and a fine-tuned encoder.}
Against the frozen gold, trained annotators reach $\kappa = 0.859$ in
leave-one-annotator-out, the best of the four LLM judges reaches $0.581$, and a
fine-tuned Legal-BERT reaches $0.720$ on \Tn{11} (Table~\ref{tab:systems}).
The judges span $\kappa$ from 0.27 to 0.58 and keep their ranking under \Tn{11},
gaining only 0.01--0.04 where the annotators recover 0.067; the two strongest are
nearly redundant ($\kappa = 0.80$ between them, above any human pair), so the
benchmark holds three distinct behaviours rather than four.
Reported rather than demonstrated (Section~\ref{sec:discussion}), their ordering
is nonetheless unambiguous: a compact, domain-specific encoder outperforms every
judge while remaining well below the human ceiling.
